
\documentclass{article}
\usepackage{colortbl}
\usepackage{makecell}
\usepackage{multirow}
\usepackage{supertabular}

\begin{document}

\newcounter{utterance}

\centering \large Interaction Transcript for game `codenames', experiment `association\_difficult', episode 1 with claude{-}opus{-}4{-}8{-}azure.
\vspace{24pt}

{ \footnotesize  \setcounter{utterance}{1}
\setlength{\tabcolsep}{0pt}
\begin{supertabular}{c@{$\;$}|p{.15\linewidth}@{}p{.15\linewidth}p{.15\linewidth}p{.15\linewidth}p{.15\linewidth}p{.15\linewidth}}
    \# & $\;$A & \multicolumn{4}{c}{Game Master} & $\;\:$B\\
    \hline

    \theutterance \stepcounter{utterance}  
    & & & \multicolumn{2}{p{0.3\linewidth}}{
        \cellcolor[rgb]{0.9,0.9,0.9}{
            \makecell[{{p{\linewidth}}}]{
                \texttt{\tiny{[GM$|$GM]}}
                {'hidden': {'team': ['zilver', 'sering', 'beha', 'tipi', 'suiker', 'verduistering', 'duplex', 'das', 'narcis'], 'innocent': ['plafond', 'tuberculose', 'pioenroos', 'grot', 'kamer', 't-shirt', 'pest'], 'opponent': ['muur', 'petunia', 'wit', 'snelheidsovertreding', 'hut', 'mosterd', 'dak', 'riem'], 'assassin': ['ontvoering']}, 'revealed': {'team': {'team': [], 'innocent': [], 'opponent': [], 'assassin': []}, 'opponent': {'team': [], 'innocent': [], 'opponent': [], 'assassin': []}}}
            }
        }
    }
    & & \\ \\

    \theutterance \stepcounter{utterance}  
    & & \multicolumn{4}{p{0.6\linewidth}}{
        \cellcolor[rgb]{0.9,0.9,0.9}{
            \makecell[{{p{\linewidth}}}]{
                \texttt{\tiny{[P1$\langle$GM]}}
                \texttt{Geef één enkele woordaanwijzing die betrekking heeft op enkele woorden van uw team uit de onderstaande lijst. U kunt kiezen op hoeveel en welke woorden uw aanwijzing betrekking heeft.} \\
\texttt{Het doel is om uw partner alle woorden van uw team te laten raden op basis van de aanwijzingen die u geeft, voordat het andere team al hun woorden heeft ontdekt. Als uw partner een 'tegenstander' of een 'afleider' raadt, eindigt hun raadbeurt en mag het andere team een beurt doen. Als uw partner een 'moord' woord raadt, eindigt het spel en verliest uw team onmiddellijk.} \\
\texttt{Hoe meer woorden uw aanwijzing met elkaar in verband brengt, hoe sneller u het spel kunt winnen, maar hoe moeilijker de associaties voor uw partner te raden zijn, dus kies uw aanwijzing verstandig. Het aanwijzingswoord moet semantisch gerelateerd zijn aan de doelwoorden, het mag niet één van de woorden in de lijsten zijn of delen ervan bevatten.} \\
\\ 
\texttt{Geef uw aanwijzing van één woord en uw door komma's gescheiden lijst van gerelateerde doelwoorden altijd in het volgende formaat en maak uw antwoorden zo kort mogelijk, voeg nooit andere tekst toe dan wat in dit formulier is vereist:} \\
\\ 
\texttt{HINT: <WOORD>} \\
\texttt{DOELWITTEN: <DOELWITTEN>} \\
\\ 
\texttt{Uw teamwoorden zijn: zilver, sering, beha, tipi, suiker, verduistering, duplex, das, narcis.} \\
\texttt{De teamwoorden van uw tegenstander zijn: muur, petunia, wit, snelheidsovertreding, hut, mosterd, dak, riem.} \\
\texttt{Afleidingswoorden zijn: plafond, tuberculose, pioenroos, grot, kamer, t{-}shirt, pest.} \\
\texttt{De moordwoorden zijn: ontvoering.} \\
            }
        }
    }
    & \\ \\

    \theutterance \stepcounter{utterance}  
    & \multicolumn{4}{p{0.6\linewidth}}{
        \cellcolor[rgb]{0.9,0.9,0.9}{
            \makecell[{{p{\linewidth}}}]{
                \texttt{\tiny{[P1$\rangle$GM]}}
                \texttt{HINT: bloesem} \\
\texttt{DOELWITTEN: sering, narcis} \\
            }
        }
    }
    & & \\ \\

    \theutterance \stepcounter{utterance}  
    & & & \multicolumn{2}{p{0.3\linewidth}}{
        \cellcolor[rgb]{0.9,0.9,0.9}{
            \makecell[{{p{\linewidth}}}]{
                \texttt{\tiny{[GM$|$GM]}}
                \texttt{bloesem} \\
            }
        }
    }
    & & \\ \\

    \theutterance \stepcounter{utterance}  
    & & & \multicolumn{2}{p{0.3\linewidth}}{
        \cellcolor[rgb]{0.9,0.9,0.9}{
            \makecell[{{p{\linewidth}}}]{
                \texttt{\tiny{[GM$|$GM]}}
                ['sering', 'narcis']
            }
        }
    }
    & & \\ \\

    \theutterance \stepcounter{utterance}  
    & & & \multicolumn{2}{p{0.3\linewidth}}{
        \cellcolor[rgb]{0.9,0.9,0.9}{
            \makecell[{{p{\linewidth}}}]{
                \texttt{\tiny{[GM$|$GM]}}
                {'word': 'sering', 'assignment': 'team'}
            }
        }
    }
    & & \\ \\

    \theutterance \stepcounter{utterance}  
    & & & \multicolumn{2}{p{0.3\linewidth}}{
        \cellcolor[rgb]{0.9,0.9,0.9}{
            \makecell[{{p{\linewidth}}}]{
                \texttt{\tiny{[GM$|$GM]}}
                {'word': 'narcis', 'assignment': 'team'}
            }
        }
    }
    & & \\ \\

    \theutterance \stepcounter{utterance}  
    & & \multicolumn{4}{p{0.6\linewidth}}{
        \cellcolor[rgb]{0.9,0.9,0.9}{
            \makecell[{{p{\linewidth}}}]{
                \texttt{\tiny{[GM$\rangle$P2]}}
                \texttt{Geef een door komma's gescheiden lijst van maximaal 2 woorden uit de volgende lijst die het beste betrekking hebben op of het meest geassocieerd zijn met het woord 'bloesem'. Begin je lijst met gissingen altijd met GISSING:  en voeg geen andere tekst toe aan je antwoord.} \\
\\ 
\texttt{grot, mosterd, riem, dak, muur, suiker, ontvoering, kamer, tipi, t{-}shirt, duplex, plafond, wit, sering, verduistering, snelheidsovertreding, tuberculose, narcis, pioenroos, zilver, petunia, pest, hut, das, beha} \\
            }
        }
    }
    & \\ \\

    \theutterance \stepcounter{utterance}  
    & & & \multicolumn{4}{p{0.6\linewidth}}{
        \cellcolor[rgb]{0.9,0.9,0.9}{
            \makecell[{{p{\linewidth}}}]{
                \texttt{\tiny{[GM$\langle$P2]}}
                \texttt{GISSING: sering, pioenroos} \\
            }
        }
    }
     \\ \\

    \theutterance \stepcounter{utterance}  
    & & & \multicolumn{2}{p{0.3\linewidth}}{
        \cellcolor[rgb]{0.9,0.9,0.9}{
            \makecell[{{p{\linewidth}}}]{
                \texttt{\tiny{[GM$|$GM]}}
                ['sering', 'pioenroos']
            }
        }
    }
    & & \\ \\

    \theutterance \stepcounter{utterance}  
    & & & \multicolumn{2}{p{0.3\linewidth}}{
        \cellcolor[rgb]{0.9,0.9,0.9}{
            \makecell[{{p{\linewidth}}}]{
                \texttt{\tiny{[GM$|$GM]}}
                {'word': 'sering', 'assignment': 'team'}
            }
        }
    }
    & & \\ \\

    \theutterance \stepcounter{utterance}  
    & & & \multicolumn{2}{p{0.3\linewidth}}{
        \cellcolor[rgb]{0.9,0.9,0.9}{
            \makecell[{{p{\linewidth}}}]{
                \texttt{\tiny{[GM$|$GM]}}
                {'word': 'sering', 'assignment': 'team'}
            }
        }
    }
    & & \\ \\

    \theutterance \stepcounter{utterance}  
    & & & \multicolumn{2}{p{0.3\linewidth}}{
        \cellcolor[rgb]{0.9,0.9,0.9}{
            \makecell[{{p{\linewidth}}}]{
                \texttt{\tiny{[GM$|$GM]}}
                {'word': 'pioenroos', 'assignment': 'innocent'}
            }
        }
    }
    & & \\ \\

    \theutterance \stepcounter{utterance}  
    & & & \multicolumn{2}{p{0.3\linewidth}}{
        \cellcolor[rgb]{0.9,0.9,0.9}{
            \makecell[{{p{\linewidth}}}]{
                \texttt{\tiny{[GM$|$GM]}}
                \texttt{pioenroos} \\
            }
        }
    }
    & & \\ \\

    \theutterance \stepcounter{utterance}  
    & & & \multicolumn{2}{p{0.3\linewidth}}{
        \cellcolor[rgb]{0.9,0.9,0.9}{
            \makecell[{{p{\linewidth}}}]{
                \texttt{\tiny{[GM$|$GM]}}
                {'word': 'snelheidsovertreding', 'assignment': 'opponent'}
            }
        }
    }
    & & \\ \\

    \theutterance \stepcounter{utterance}  
    & & & \multicolumn{2}{p{0.3\linewidth}}{
        \cellcolor[rgb]{0.9,0.9,0.9}{
            \makecell[{{p{\linewidth}}}]{
                \texttt{\tiny{[GM$|$GM]}}
                {'hidden': {'team': ['zilver', 'beha', 'tipi', 'suiker', 'verduistering', 'duplex', 'das', 'narcis'], 'innocent': ['plafond', 'tuberculose', 'grot', 'kamer', 't-shirt', 'pest'], 'opponent': ['muur', 'petunia', 'wit', 'hut', 'mosterd', 'dak', 'riem'], 'assassin': ['ontvoering']}, 'revealed': {'team': {'team': ['sering'], 'innocent': ['pioenroos'], 'opponent': [], 'assassin': []}, 'opponent': {'team': [], 'innocent': [], 'opponent': ['snelheidsovertreding'], 'assassin': []}}}
            }
        }
    }
    & & \\ \\

    \theutterance \stepcounter{utterance}  
    & & \multicolumn{4}{p{0.6\linewidth}}{
        \cellcolor[rgb]{0.9,0.9,0.9}{
            \makecell[{{p{\linewidth}}}]{
                \texttt{\tiny{[P1$\langle$GM]}}
                \texttt{Het woord sering is correct geraden. Het woord pioenroos is geraden, maar het is een onschuldig woord. De beurt van je teamgenoot is daar geëindigd.} \\
\texttt{Geef nu nog een aanwijzing die betrekking heeft op enkele van je overgebleven teamwoorden en een lijst met de bijbehorende doelwoorden. Vergeet niet je aanwijzing te beginnen met 'HINT: ', zet een nieuwe regel en begin je door komma's gescheiden lijst van doelwoorden met 'DOELWITTEN: '. Let op: sommige woorden zijn verwijderd uit de lijsten vergeleken met eerdere verzoeken.} \\
\\ 
\texttt{Je resterende teamwoorden zijn: zilver, beha, tipi, suiker, verduistering, duplex, das, narcis.} \\
\texttt{Resterende woorden voor je tegenstander zijn: muur, petunia, wit, snelheidsovertreding, hut, mosterd, dak, riem.} \\
\texttt{Resterende afleidingswoorden zijn: plafond, tuberculose, grot, kamer, t{-}shirt, pest.} \\
\texttt{Resterende killwoorden zijn: ontvoering.} \\
            }
        }
    }
    & \\ \\

    \theutterance \stepcounter{utterance}  
    & \multicolumn{4}{p{0.6\linewidth}}{
        \cellcolor[rgb]{0.9,0.9,0.9}{
            \makecell[{{p{\linewidth}}}]{
                \texttt{\tiny{[P1$\rangle$GM]}}
                \texttt{HINT: maan} \\
\texttt{DOELWITTEN: zilver, verduistering} \\
            }
        }
    }
    & & \\ \\

    \theutterance \stepcounter{utterance}  
    & & & \multicolumn{2}{p{0.3\linewidth}}{
        \cellcolor[rgb]{0.9,0.9,0.9}{
            \makecell[{{p{\linewidth}}}]{
                \texttt{\tiny{[GM$|$GM]}}
                \texttt{maan} \\
            }
        }
    }
    & & \\ \\

    \theutterance \stepcounter{utterance}  
    & & & \multicolumn{2}{p{0.3\linewidth}}{
        \cellcolor[rgb]{0.9,0.9,0.9}{
            \makecell[{{p{\linewidth}}}]{
                \texttt{\tiny{[GM$|$GM]}}
                ['zilver', 'verduistering']
            }
        }
    }
    & & \\ \\

    \theutterance \stepcounter{utterance}  
    & & & \multicolumn{2}{p{0.3\linewidth}}{
        \cellcolor[rgb]{0.9,0.9,0.9}{
            \makecell[{{p{\linewidth}}}]{
                \texttt{\tiny{[GM$|$GM]}}
                {'word': 'zilver', 'assignment': 'team'}
            }
        }
    }
    & & \\ \\

    \theutterance \stepcounter{utterance}  
    & & & \multicolumn{2}{p{0.3\linewidth}}{
        \cellcolor[rgb]{0.9,0.9,0.9}{
            \makecell[{{p{\linewidth}}}]{
                \texttt{\tiny{[GM$|$GM]}}
                {'word': 'verduistering', 'assignment': 'team'}
            }
        }
    }
    & & \\ \\

    \theutterance \stepcounter{utterance}  
    & & \multicolumn{4}{p{0.6\linewidth}}{
        \cellcolor[rgb]{0.9,0.9,0.9}{
            \makecell[{{p{\linewidth}}}]{
                \texttt{\tiny{[GM$\rangle$P2]}}
                \texttt{Het woord sering is correct geraden. Het woord pioenroos is geraden, maar het is een onschuldig woord. Je beurt is daar geëindigd.} \\
\texttt{Geef nu nog een door komma's gescheiden lijst van ten minste 1 en maximaal 2 woorden uit de volgende lijst van woorden die het beste betrekking hebben op of het meest geassocieerd zijn met het woord 'maan'. Vergeet niet uw antwoord te beginnen met GISSING: . Opmerking: sommige woorden zijn verwijderd uit de lijst vergeleken met eerdere verzoeken.} \\
\\ 
\texttt{grot, mosterd, riem, dak, muur, suiker, ontvoering, kamer, tipi, t{-}shirt, duplex, plafond, wit, verduistering, tuberculose, narcis, zilver, petunia, pest, hut, das, beha} \\
            }
        }
    }
    & \\ \\

    \theutterance \stepcounter{utterance}  
    & & & \multicolumn{4}{p{0.6\linewidth}}{
        \cellcolor[rgb]{0.9,0.9,0.9}{
            \makecell[{{p{\linewidth}}}]{
                \texttt{\tiny{[GM$\langle$P2]}}
                \texttt{GISSING: verduistering, zilver} \\
            }
        }
    }
     \\ \\

    \theutterance \stepcounter{utterance}  
    & & & \multicolumn{2}{p{0.3\linewidth}}{
        \cellcolor[rgb]{0.9,0.9,0.9}{
            \makecell[{{p{\linewidth}}}]{
                \texttt{\tiny{[GM$|$GM]}}
                ['verduistering', 'zilver']
            }
        }
    }
    & & \\ \\

    \theutterance \stepcounter{utterance}  
    & & & \multicolumn{2}{p{0.3\linewidth}}{
        \cellcolor[rgb]{0.9,0.9,0.9}{
            \makecell[{{p{\linewidth}}}]{
                \texttt{\tiny{[GM$|$GM]}}
                {'word': 'verduistering', 'assignment': 'team'}
            }
        }
    }
    & & \\ \\

    \theutterance \stepcounter{utterance}  
    & & & \multicolumn{2}{p{0.3\linewidth}}{
        \cellcolor[rgb]{0.9,0.9,0.9}{
            \makecell[{{p{\linewidth}}}]{
                \texttt{\tiny{[GM$|$GM]}}
                {'word': 'verduistering', 'assignment': 'team'}
            }
        }
    }
    & & \\ \\

    \theutterance \stepcounter{utterance}  
    & & & \multicolumn{2}{p{0.3\linewidth}}{
        \cellcolor[rgb]{0.9,0.9,0.9}{
            \makecell[{{p{\linewidth}}}]{
                \texttt{\tiny{[GM$|$GM]}}
                {'word': 'zilver', 'assignment': 'team'}
            }
        }
    }
    & & \\ \\

    \theutterance \stepcounter{utterance}  
    & & & \multicolumn{2}{p{0.3\linewidth}}{
        \cellcolor[rgb]{0.9,0.9,0.9}{
            \makecell[{{p{\linewidth}}}]{
                \texttt{\tiny{[GM$|$GM]}}
                {'word': 'zilver', 'assignment': 'team'}
            }
        }
    }
    & & \\ \\

    \theutterance \stepcounter{utterance}  
    & & & \multicolumn{2}{p{0.3\linewidth}}{
        \cellcolor[rgb]{0.9,0.9,0.9}{
            \makecell[{{p{\linewidth}}}]{
                \texttt{\tiny{[GM$|$GM]}}
                {'word': 'riem', 'assignment': 'opponent'}
            }
        }
    }
    & & \\ \\

    \theutterance \stepcounter{utterance}  
    & & & \multicolumn{2}{p{0.3\linewidth}}{
        \cellcolor[rgb]{0.9,0.9,0.9}{
            \makecell[{{p{\linewidth}}}]{
                \texttt{\tiny{[GM$|$GM]}}
                {'hidden': {'team': ['beha', 'tipi', 'suiker', 'duplex', 'das', 'narcis'], 'innocent': ['plafond', 'tuberculose', 'grot', 'kamer', 't-shirt', 'pest'], 'opponent': ['muur', 'petunia', 'wit', 'hut', 'mosterd', 'dak'], 'assassin': ['ontvoering']}, 'revealed': {'team': {'team': ['sering', 'verduistering', 'zilver'], 'innocent': ['pioenroos'], 'opponent': [], 'assassin': []}, 'opponent': {'team': [], 'innocent': [], 'opponent': ['snelheidsovertreding', 'riem'], 'assassin': []}}}
            }
        }
    }
    & & \\ \\

    \theutterance \stepcounter{utterance}  
    & & \multicolumn{4}{p{0.6\linewidth}}{
        \cellcolor[rgb]{0.9,0.9,0.9}{
            \makecell[{{p{\linewidth}}}]{
                \texttt{\tiny{[P1$\langle$GM]}}
                \texttt{De woorden verduistering, zilver zijn correct geraden. De beurt van je teamgenoot is daar geëindigd.} \\
\texttt{Geef nu nog een aanwijzing die betrekking heeft op enkele van je overgebleven teamwoorden en een lijst met de bijbehorende doelwoorden. Vergeet niet je aanwijzing te beginnen met 'HINT: ', zet een nieuwe regel en begin je door komma's gescheiden lijst van doelwoorden met 'DOELWITTEN: '. Let op: sommige woorden zijn verwijderd uit de lijsten vergeleken met eerdere verzoeken.} \\
\\ 
\texttt{Je resterende teamwoorden zijn: beha, tipi, suiker, duplex, das, narcis.} \\
\texttt{Resterende woorden voor je tegenstander zijn: muur, petunia, wit, hut, mosterd, dak, riem.} \\
\texttt{Resterende afleidingswoorden zijn: plafond, tuberculose, grot, kamer, t{-}shirt, pest.} \\
\texttt{Resterende killwoorden zijn: ontvoering.} \\
            }
        }
    }
    & \\ \\

    \theutterance \stepcounter{utterance}  
    & \multicolumn{4}{p{0.6\linewidth}}{
        \cellcolor[rgb]{0.9,0.9,0.9}{
            \makecell[{{p{\linewidth}}}]{
                \texttt{\tiny{[P1$\rangle$GM]}}
                \texttt{HINT: kleding} \\
\texttt{DOELWITTEN: beha, das} \\
            }
        }
    }
    & & \\ \\

    \theutterance \stepcounter{utterance}  
    & & & \multicolumn{2}{p{0.3\linewidth}}{
        \cellcolor[rgb]{0.9,0.9,0.9}{
            \makecell[{{p{\linewidth}}}]{
                \texttt{\tiny{[GM$|$GM]}}
                \texttt{kleding} \\
            }
        }
    }
    & & \\ \\

    \theutterance \stepcounter{utterance}  
    & & & \multicolumn{2}{p{0.3\linewidth}}{
        \cellcolor[rgb]{0.9,0.9,0.9}{
            \makecell[{{p{\linewidth}}}]{
                \texttt{\tiny{[GM$|$GM]}}
                ['beha', 'das']
            }
        }
    }
    & & \\ \\

    \theutterance \stepcounter{utterance}  
    & & & \multicolumn{2}{p{0.3\linewidth}}{
        \cellcolor[rgb]{0.9,0.9,0.9}{
            \makecell[{{p{\linewidth}}}]{
                \texttt{\tiny{[GM$|$GM]}}
                {'word': 'beha', 'assignment': 'team'}
            }
        }
    }
    & & \\ \\

    \theutterance \stepcounter{utterance}  
    & & & \multicolumn{2}{p{0.3\linewidth}}{
        \cellcolor[rgb]{0.9,0.9,0.9}{
            \makecell[{{p{\linewidth}}}]{
                \texttt{\tiny{[GM$|$GM]}}
                {'word': 'das', 'assignment': 'team'}
            }
        }
    }
    & & \\ \\

    \theutterance \stepcounter{utterance}  
    & & \multicolumn{4}{p{0.6\linewidth}}{
        \cellcolor[rgb]{0.9,0.9,0.9}{
            \makecell[{{p{\linewidth}}}]{
                \texttt{\tiny{[GM$\rangle$P2]}}
                \texttt{De woorden verduistering, zilver zijn correct geraden. Je beurt is daar geëindigd.} \\
\texttt{Geef nu nog een door komma's gescheiden lijst van ten minste 1 en maximaal 2 woorden uit de volgende lijst van woorden die het beste betrekking hebben op of het meest geassocieerd zijn met het woord 'kleding'. Vergeet niet uw antwoord te beginnen met GISSING: . Opmerking: sommige woorden zijn verwijderd uit de lijst vergeleken met eerdere verzoeken.} \\
\\ 
\texttt{grot, mosterd, dak, muur, suiker, ontvoering, kamer, tipi, t{-}shirt, duplex, plafond, wit, tuberculose, narcis, petunia, pest, hut, das, beha} \\
            }
        }
    }
    & \\ \\

    \theutterance \stepcounter{utterance}  
    & & & \multicolumn{4}{p{0.6\linewidth}}{
        \cellcolor[rgb]{0.9,0.9,0.9}{
            \makecell[{{p{\linewidth}}}]{
                \texttt{\tiny{[GM$\langle$P2]}}
                \texttt{GISSING: t{-}shirt, beha} \\
            }
        }
    }
     \\ \\

    \theutterance \stepcounter{utterance}  
    & & & \multicolumn{2}{p{0.3\linewidth}}{
        \cellcolor[rgb]{0.9,0.9,0.9}{
            \makecell[{{p{\linewidth}}}]{
                \texttt{\tiny{[GM$|$GM]}}
                ['t-shirt', 'beha']
            }
        }
    }
    & & \\ \\

    \theutterance \stepcounter{utterance}  
    & & & \multicolumn{2}{p{0.3\linewidth}}{
        \cellcolor[rgb]{0.9,0.9,0.9}{
            \makecell[{{p{\linewidth}}}]{
                \texttt{\tiny{[GM$|$GM]}}
                {'word': 't-shirt', 'assignment': 'innocent'}
            }
        }
    }
    & & \\ \\

    \theutterance \stepcounter{utterance}  
    & & & \multicolumn{2}{p{0.3\linewidth}}{
        \cellcolor[rgb]{0.9,0.9,0.9}{
            \makecell[{{p{\linewidth}}}]{
                \texttt{\tiny{[GM$|$GM]}}
                \texttt{t{-}shirt} \\
            }
        }
    }
    & & \\ \\

    \theutterance \stepcounter{utterance}  
    & & & \multicolumn{2}{p{0.3\linewidth}}{
        \cellcolor[rgb]{0.9,0.9,0.9}{
            \makecell[{{p{\linewidth}}}]{
                \texttt{\tiny{[GM$|$GM]}}
                {'word': 'wit', 'assignment': 'opponent'}
            }
        }
    }
    & & \\ \\

    \theutterance \stepcounter{utterance}  
    & & & \multicolumn{2}{p{0.3\linewidth}}{
        \cellcolor[rgb]{0.9,0.9,0.9}{
            \makecell[{{p{\linewidth}}}]{
                \texttt{\tiny{[GM$|$GM]}}
                {'hidden': {'team': ['beha', 'tipi', 'suiker', 'duplex', 'das', 'narcis'], 'innocent': ['plafond', 'tuberculose', 'grot', 'kamer', 'pest'], 'opponent': ['muur', 'petunia', 'hut', 'mosterd', 'dak'], 'assassin': ['ontvoering']}, 'revealed': {'team': {'team': ['sering', 'verduistering', 'zilver'], 'innocent': ['pioenroos', 't-shirt'], 'opponent': [], 'assassin': []}, 'opponent': {'team': [], 'innocent': [], 'opponent': ['snelheidsovertreding', 'riem', 'wit'], 'assassin': []}}}
            }
        }
    }
    & & \\ \\

    \theutterance \stepcounter{utterance}  
    & & \multicolumn{4}{p{0.6\linewidth}}{
        \cellcolor[rgb]{0.9,0.9,0.9}{
            \makecell[{{p{\linewidth}}}]{
                \texttt{\tiny{[P1$\langle$GM]}}
                \texttt{Het woord t{-}shirt is geraden, maar het is een onschuldig woord. De beurt van je teamgenoot is daar geëindigd.} \\
\texttt{Geef nu nog een aanwijzing die betrekking heeft op enkele van je overgebleven teamwoorden en een lijst met de bijbehorende doelwoorden. Vergeet niet je aanwijzing te beginnen met 'HINT: ', zet een nieuwe regel en begin je door komma's gescheiden lijst van doelwoorden met 'DOELWITTEN: '. Let op: sommige woorden zijn verwijderd uit de lijsten vergeleken met eerdere verzoeken.} \\
\\ 
\texttt{Je resterende teamwoorden zijn: beha, tipi, suiker, duplex, das, narcis.} \\
\texttt{Resterende woorden voor je tegenstander zijn: muur, petunia, wit, hut, mosterd, dak.} \\
\texttt{Resterende afleidingswoorden zijn: plafond, tuberculose, grot, kamer, pest.} \\
\texttt{Resterende killwoorden zijn: ontvoering.} \\
            }
        }
    }
    & \\ \\

    \theutterance \stepcounter{utterance}  
    & \multicolumn{4}{p{0.6\linewidth}}{
        \cellcolor[rgb]{0.9,0.9,0.9}{
            \makecell[{{p{\linewidth}}}]{
                \texttt{\tiny{[P1$\rangle$GM]}}
                \texttt{HINT: woning} \\
\texttt{DOELWITTEN: tipi, duplex} \\
            }
        }
    }
    & & \\ \\

    \theutterance \stepcounter{utterance}  
    & & & \multicolumn{2}{p{0.3\linewidth}}{
        \cellcolor[rgb]{0.9,0.9,0.9}{
            \makecell[{{p{\linewidth}}}]{
                \texttt{\tiny{[GM$|$GM]}}
                \texttt{woning} \\
            }
        }
    }
    & & \\ \\

    \theutterance \stepcounter{utterance}  
    & & & \multicolumn{2}{p{0.3\linewidth}}{
        \cellcolor[rgb]{0.9,0.9,0.9}{
            \makecell[{{p{\linewidth}}}]{
                \texttt{\tiny{[GM$|$GM]}}
                ['tipi', 'duplex']
            }
        }
    }
    & & \\ \\

    \theutterance \stepcounter{utterance}  
    & & & \multicolumn{2}{p{0.3\linewidth}}{
        \cellcolor[rgb]{0.9,0.9,0.9}{
            \makecell[{{p{\linewidth}}}]{
                \texttt{\tiny{[GM$|$GM]}}
                {'word': 'tipi', 'assignment': 'team'}
            }
        }
    }
    & & \\ \\

    \theutterance \stepcounter{utterance}  
    & & & \multicolumn{2}{p{0.3\linewidth}}{
        \cellcolor[rgb]{0.9,0.9,0.9}{
            \makecell[{{p{\linewidth}}}]{
                \texttt{\tiny{[GM$|$GM]}}
                {'word': 'duplex', 'assignment': 'team'}
            }
        }
    }
    & & \\ \\

    \theutterance \stepcounter{utterance}  
    & & \multicolumn{4}{p{0.6\linewidth}}{
        \cellcolor[rgb]{0.9,0.9,0.9}{
            \makecell[{{p{\linewidth}}}]{
                \texttt{\tiny{[GM$\rangle$P2]}}
                \texttt{Het woord t{-}shirt is geraden, maar het is een onschuldig woord. Je beurt is daar geëindigd.} \\
\texttt{Geef nu nog een door komma's gescheiden lijst van ten minste 1 en maximaal 2 woorden uit de volgende lijst van woorden die het beste betrekking hebben op of het meest geassocieerd zijn met het woord 'woning'. Vergeet niet uw antwoord te beginnen met GISSING: . Opmerking: sommige woorden zijn verwijderd uit de lijst vergeleken met eerdere verzoeken.} \\
\\ 
\texttt{grot, mosterd, dak, muur, suiker, ontvoering, kamer, tipi, duplex, plafond, tuberculose, narcis, petunia, pest, hut, das, beha} \\
            }
        }
    }
    & \\ \\

    \theutterance \stepcounter{utterance}  
    & & & \multicolumn{4}{p{0.6\linewidth}}{
        \cellcolor[rgb]{0.9,0.9,0.9}{
            \makecell[{{p{\linewidth}}}]{
                \texttt{\tiny{[GM$\langle$P2]}}
                \texttt{GISSING: duplex, hut} \\
            }
        }
    }
     \\ \\

    \theutterance \stepcounter{utterance}  
    & & & \multicolumn{2}{p{0.3\linewidth}}{
        \cellcolor[rgb]{0.9,0.9,0.9}{
            \makecell[{{p{\linewidth}}}]{
                \texttt{\tiny{[GM$|$GM]}}
                ['duplex', 'hut']
            }
        }
    }
    & & \\ \\

    \theutterance \stepcounter{utterance}  
    & & & \multicolumn{2}{p{0.3\linewidth}}{
        \cellcolor[rgb]{0.9,0.9,0.9}{
            \makecell[{{p{\linewidth}}}]{
                \texttt{\tiny{[GM$|$GM]}}
                {'word': 'duplex', 'assignment': 'team'}
            }
        }
    }
    & & \\ \\

    \theutterance \stepcounter{utterance}  
    & & & \multicolumn{2}{p{0.3\linewidth}}{
        \cellcolor[rgb]{0.9,0.9,0.9}{
            \makecell[{{p{\linewidth}}}]{
                \texttt{\tiny{[GM$|$GM]}}
                {'word': 'duplex', 'assignment': 'team'}
            }
        }
    }
    & & \\ \\

    \theutterance \stepcounter{utterance}  
    & & & \multicolumn{2}{p{0.3\linewidth}}{
        \cellcolor[rgb]{0.9,0.9,0.9}{
            \makecell[{{p{\linewidth}}}]{
                \texttt{\tiny{[GM$|$GM]}}
                {'word': 'hut', 'assignment': 'opponent'}
            }
        }
    }
    & & \\ \\

    \theutterance \stepcounter{utterance}  
    & & & \multicolumn{2}{p{0.3\linewidth}}{
        \cellcolor[rgb]{0.9,0.9,0.9}{
            \makecell[{{p{\linewidth}}}]{
                \texttt{\tiny{[GM$|$GM]}}
                \texttt{hut} \\
            }
        }
    }
    & & \\ \\

    \theutterance \stepcounter{utterance}  
    & & & \multicolumn{2}{p{0.3\linewidth}}{
        \cellcolor[rgb]{0.9,0.9,0.9}{
            \makecell[{{p{\linewidth}}}]{
                \texttt{\tiny{[GM$|$GM]}}
                {'word': 'muur', 'assignment': 'opponent'}
            }
        }
    }
    & & \\ \\

    \theutterance \stepcounter{utterance}  
    & & & \multicolumn{2}{p{0.3\linewidth}}{
        \cellcolor[rgb]{0.9,0.9,0.9}{
            \makecell[{{p{\linewidth}}}]{
                \texttt{\tiny{[GM$|$GM]}}
                {'hidden': {'team': ['beha', 'tipi', 'suiker', 'das', 'narcis'], 'innocent': ['plafond', 'tuberculose', 'grot', 'kamer', 'pest'], 'opponent': ['petunia', 'mosterd', 'dak'], 'assassin': ['ontvoering']}, 'revealed': {'team': {'team': ['sering', 'verduistering', 'zilver', 'duplex'], 'innocent': ['pioenroos', 't-shirt'], 'opponent': ['hut'], 'assassin': []}, 'opponent': {'team': [], 'innocent': [], 'opponent': ['snelheidsovertreding', 'riem', 'wit', 'muur'], 'assassin': []}}}
            }
        }
    }
    & & \\ \\

    \theutterance \stepcounter{utterance}  
    & & \multicolumn{4}{p{0.6\linewidth}}{
        \cellcolor[rgb]{0.9,0.9,0.9}{
            \makecell[{{p{\linewidth}}}]{
                \texttt{\tiny{[P1$\langle$GM]}}
                \texttt{Het woord duplex is correct geraden. Het woord hut is geraden, maar het is een tegenstander woord. De beurt van je teamgenoot is daar geëindigd.} \\
\texttt{Geef nu nog een aanwijzing die betrekking heeft op enkele van je overgebleven teamwoorden en een lijst met de bijbehorende doelwoorden. Vergeet niet je aanwijzing te beginnen met 'HINT: ', zet een nieuwe regel en begin je door komma's gescheiden lijst van doelwoorden met 'DOELWITTEN: '. Let op: sommige woorden zijn verwijderd uit de lijsten vergeleken met eerdere verzoeken.} \\
\\ 
\texttt{Je resterende teamwoorden zijn: beha, tipi, suiker, das, narcis.} \\
\texttt{Resterende woorden voor je tegenstander zijn: muur, petunia, mosterd, dak.} \\
\texttt{Resterende afleidingswoorden zijn: plafond, tuberculose, grot, kamer, pest.} \\
\texttt{Resterende killwoorden zijn: ontvoering.} \\
            }
        }
    }
    & \\ \\

    \theutterance \stepcounter{utterance}  
    & \multicolumn{4}{p{0.6\linewidth}}{
        \cellcolor[rgb]{0.9,0.9,0.9}{
            \makecell[{{p{\linewidth}}}]{
                \texttt{\tiny{[P1$\rangle$GM]}}
                \texttt{HINT: indiaan} \\
\texttt{DOELWITTEN: tipi} \\
            }
        }
    }
    & & \\ \\

    \theutterance \stepcounter{utterance}  
    & & & \multicolumn{2}{p{0.3\linewidth}}{
        \cellcolor[rgb]{0.9,0.9,0.9}{
            \makecell[{{p{\linewidth}}}]{
                \texttt{\tiny{[GM$|$GM]}}
                \texttt{indiaan} \\
            }
        }
    }
    & & \\ \\

    \theutterance \stepcounter{utterance}  
    & & & \multicolumn{2}{p{0.3\linewidth}}{
        \cellcolor[rgb]{0.9,0.9,0.9}{
            \makecell[{{p{\linewidth}}}]{
                \texttt{\tiny{[GM$|$GM]}}
                ['tipi']
            }
        }
    }
    & & \\ \\

    \theutterance \stepcounter{utterance}  
    & & & \multicolumn{2}{p{0.3\linewidth}}{
        \cellcolor[rgb]{0.9,0.9,0.9}{
            \makecell[{{p{\linewidth}}}]{
                \texttt{\tiny{[GM$|$GM]}}
                {'word': 'tipi', 'assignment': 'team'}
            }
        }
    }
    & & \\ \\

    \theutterance \stepcounter{utterance}  
    & & \multicolumn{4}{p{0.6\linewidth}}{
        \cellcolor[rgb]{0.9,0.9,0.9}{
            \makecell[{{p{\linewidth}}}]{
                \texttt{\tiny{[GM$\rangle$P2]}}
                \texttt{Het woord duplex is correct geraden. Het woord hut is geraden, maar het is een tegenstander woord. Je beurt is daar geëindigd.} \\
\texttt{Geef nu nog een door komma's gescheiden lijst van ten minste 1 en maximaal 1 woorden uit de volgende lijst van woorden die het beste betrekking hebben op of het meest geassocieerd zijn met het woord 'indiaan'. Vergeet niet uw antwoord te beginnen met GISSING: . Opmerking: sommige woorden zijn verwijderd uit de lijst vergeleken met eerdere verzoeken.} \\
\\ 
\texttt{grot, mosterd, dak, suiker, ontvoering, kamer, tipi, plafond, tuberculose, narcis, petunia, pest, das, beha} \\
            }
        }
    }
    & \\ \\

    \theutterance \stepcounter{utterance}  
    & & & \multicolumn{4}{p{0.6\linewidth}}{
        \cellcolor[rgb]{0.9,0.9,0.9}{
            \makecell[{{p{\linewidth}}}]{
                \texttt{\tiny{[GM$\langle$P2]}}
                \texttt{GISSING: tipi} \\
            }
        }
    }
     \\ \\

    \theutterance \stepcounter{utterance}  
    & & & \multicolumn{2}{p{0.3\linewidth}}{
        \cellcolor[rgb]{0.9,0.9,0.9}{
            \makecell[{{p{\linewidth}}}]{
                \texttt{\tiny{[GM$|$GM]}}
                ['tipi']
            }
        }
    }
    & & \\ \\

    \theutterance \stepcounter{utterance}  
    & & & \multicolumn{2}{p{0.3\linewidth}}{
        \cellcolor[rgb]{0.9,0.9,0.9}{
            \makecell[{{p{\linewidth}}}]{
                \texttt{\tiny{[GM$|$GM]}}
                {'word': 'tipi', 'assignment': 'team'}
            }
        }
    }
    & & \\ \\

    \theutterance \stepcounter{utterance}  
    & & & \multicolumn{2}{p{0.3\linewidth}}{
        \cellcolor[rgb]{0.9,0.9,0.9}{
            \makecell[{{p{\linewidth}}}]{
                \texttt{\tiny{[GM$|$GM]}}
                {'word': 'tipi', 'assignment': 'team'}
            }
        }
    }
    & & \\ \\

    \theutterance \stepcounter{utterance}  
    & & & \multicolumn{2}{p{0.3\linewidth}}{
        \cellcolor[rgb]{0.9,0.9,0.9}{
            \makecell[{{p{\linewidth}}}]{
                \texttt{\tiny{[GM$|$GM]}}
                {'word': 'petunia', 'assignment': 'opponent'}
            }
        }
    }
    & & \\ \\

    \theutterance \stepcounter{utterance}  
    & & & \multicolumn{2}{p{0.3\linewidth}}{
        \cellcolor[rgb]{0.9,0.9,0.9}{
            \makecell[{{p{\linewidth}}}]{
                \texttt{\tiny{[GM$|$GM]}}
                {'hidden': {'team': ['beha', 'suiker', 'das', 'narcis'], 'innocent': ['plafond', 'tuberculose', 'grot', 'kamer', 'pest'], 'opponent': ['mosterd', 'dak'], 'assassin': ['ontvoering']}, 'revealed': {'team': {'team': ['sering', 'verduistering', 'zilver', 'duplex', 'tipi'], 'innocent': ['pioenroos', 't-shirt'], 'opponent': ['hut'], 'assassin': []}, 'opponent': {'team': [], 'innocent': [], 'opponent': ['snelheidsovertreding', 'riem', 'wit', 'muur', 'petunia'], 'assassin': []}}}
            }
        }
    }
    & & \\ \\

    \theutterance \stepcounter{utterance}  
    & & \multicolumn{4}{p{0.6\linewidth}}{
        \cellcolor[rgb]{0.9,0.9,0.9}{
            \makecell[{{p{\linewidth}}}]{
                \texttt{\tiny{[P1$\langle$GM]}}
                \texttt{Het woord tipi is correct geraden. De beurt van je teamgenoot is daar geëindigd.} \\
\texttt{Geef nu nog een aanwijzing die betrekking heeft op enkele van je overgebleven teamwoorden en een lijst met de bijbehorende doelwoorden. Vergeet niet je aanwijzing te beginnen met 'HINT: ', zet een nieuwe regel en begin je door komma's gescheiden lijst van doelwoorden met 'DOELWITTEN: '. Let op: sommige woorden zijn verwijderd uit de lijsten vergeleken met eerdere verzoeken.} \\
\\ 
\texttt{Je resterende teamwoorden zijn: beha, suiker, das, narcis.} \\
\texttt{Resterende woorden voor je tegenstander zijn: petunia, mosterd, dak.} \\
\texttt{Resterende afleidingswoorden zijn: plafond, tuberculose, grot, kamer, pest.} \\
\texttt{Resterende killwoorden zijn: ontvoering.} \\
            }
        }
    }
    & \\ \\

    \theutterance \stepcounter{utterance}  
    & \multicolumn{4}{p{0.6\linewidth}}{
        \cellcolor[rgb]{0.9,0.9,0.9}{
            \makecell[{{p{\linewidth}}}]{
                \texttt{\tiny{[P1$\rangle$GM]}}
                \texttt{HINT: zoet} \\
\texttt{DOELWITTEN: suiker} \\
            }
        }
    }
    & & \\ \\

    \theutterance \stepcounter{utterance}  
    & & & \multicolumn{2}{p{0.3\linewidth}}{
        \cellcolor[rgb]{0.9,0.9,0.9}{
            \makecell[{{p{\linewidth}}}]{
                \texttt{\tiny{[GM$|$GM]}}
                \texttt{zoet} \\
            }
        }
    }
    & & \\ \\

    \theutterance \stepcounter{utterance}  
    & & & \multicolumn{2}{p{0.3\linewidth}}{
        \cellcolor[rgb]{0.9,0.9,0.9}{
            \makecell[{{p{\linewidth}}}]{
                \texttt{\tiny{[GM$|$GM]}}
                ['suiker']
            }
        }
    }
    & & \\ \\

    \theutterance \stepcounter{utterance}  
    & & & \multicolumn{2}{p{0.3\linewidth}}{
        \cellcolor[rgb]{0.9,0.9,0.9}{
            \makecell[{{p{\linewidth}}}]{
                \texttt{\tiny{[GM$|$GM]}}
                {'word': 'suiker', 'assignment': 'team'}
            }
        }
    }
    & & \\ \\

    \theutterance \stepcounter{utterance}  
    & & \multicolumn{4}{p{0.6\linewidth}}{
        \cellcolor[rgb]{0.9,0.9,0.9}{
            \makecell[{{p{\linewidth}}}]{
                \texttt{\tiny{[GM$\rangle$P2]}}
                \texttt{Het woord tipi is correct geraden. Je beurt is daar geëindigd.} \\
\texttt{Geef nu nog een door komma's gescheiden lijst van ten minste 1 en maximaal 1 woorden uit de volgende lijst van woorden die het beste betrekking hebben op of het meest geassocieerd zijn met het woord 'zoet'. Vergeet niet uw antwoord te beginnen met GISSING: . Opmerking: sommige woorden zijn verwijderd uit de lijst vergeleken met eerdere verzoeken.} \\
\\ 
\texttt{grot, mosterd, dak, suiker, ontvoering, kamer, plafond, tuberculose, narcis, pest, das, beha} \\
            }
        }
    }
    & \\ \\

    \theutterance \stepcounter{utterance}  
    & & & \multicolumn{4}{p{0.6\linewidth}}{
        \cellcolor[rgb]{0.9,0.9,0.9}{
            \makecell[{{p{\linewidth}}}]{
                \texttt{\tiny{[GM$\langle$P2]}}
                \texttt{GISSING: suiker} \\
            }
        }
    }
     \\ \\

    \theutterance \stepcounter{utterance}  
    & & & \multicolumn{2}{p{0.3\linewidth}}{
        \cellcolor[rgb]{0.9,0.9,0.9}{
            \makecell[{{p{\linewidth}}}]{
                \texttt{\tiny{[GM$|$GM]}}
                ['suiker']
            }
        }
    }
    & & \\ \\

    \theutterance \stepcounter{utterance}  
    & & & \multicolumn{2}{p{0.3\linewidth}}{
        \cellcolor[rgb]{0.9,0.9,0.9}{
            \makecell[{{p{\linewidth}}}]{
                \texttt{\tiny{[GM$|$GM]}}
                {'word': 'suiker', 'assignment': 'team'}
            }
        }
    }
    & & \\ \\

    \theutterance \stepcounter{utterance}  
    & & & \multicolumn{2}{p{0.3\linewidth}}{
        \cellcolor[rgb]{0.9,0.9,0.9}{
            \makecell[{{p{\linewidth}}}]{
                \texttt{\tiny{[GM$|$GM]}}
                {'word': 'suiker', 'assignment': 'team'}
            }
        }
    }
    & & \\ \\

    \theutterance \stepcounter{utterance}  
    & & & \multicolumn{2}{p{0.3\linewidth}}{
        \cellcolor[rgb]{0.9,0.9,0.9}{
            \makecell[{{p{\linewidth}}}]{
                \texttt{\tiny{[GM$|$GM]}}
                {'word': 'mosterd', 'assignment': 'opponent'}
            }
        }
    }
    & & \\ \\

    \theutterance \stepcounter{utterance}  
    & & & \multicolumn{2}{p{0.3\linewidth}}{
        \cellcolor[rgb]{0.9,0.9,0.9}{
            \makecell[{{p{\linewidth}}}]{
                \texttt{\tiny{[GM$|$GM]}}
                {'hidden': {'team': ['beha', 'das', 'narcis'], 'innocent': ['plafond', 'tuberculose', 'grot', 'kamer', 'pest'], 'opponent': ['dak'], 'assassin': ['ontvoering']}, 'revealed': {'team': {'team': ['sering', 'verduistering', 'zilver', 'duplex', 'tipi', 'suiker'], 'innocent': ['pioenroos', 't-shirt'], 'opponent': ['hut'], 'assassin': []}, 'opponent': {'team': [], 'innocent': [], 'opponent': ['snelheidsovertreding', 'riem', 'wit', 'muur', 'petunia', 'mosterd'], 'assassin': []}}}
            }
        }
    }
    & & \\ \\

    \theutterance \stepcounter{utterance}  
    & & \multicolumn{4}{p{0.6\linewidth}}{
        \cellcolor[rgb]{0.9,0.9,0.9}{
            \makecell[{{p{\linewidth}}}]{
                \texttt{\tiny{[P1$\langle$GM]}}
                \texttt{Het woord suiker is correct geraden. De beurt van je teamgenoot is daar geëindigd.} \\
\texttt{Geef nu nog een aanwijzing die betrekking heeft op enkele van je overgebleven teamwoorden en een lijst met de bijbehorende doelwoorden. Vergeet niet je aanwijzing te beginnen met 'HINT: ', zet een nieuwe regel en begin je door komma's gescheiden lijst van doelwoorden met 'DOELWITTEN: '. Let op: sommige woorden zijn verwijderd uit de lijsten vergeleken met eerdere verzoeken.} \\
\\ 
\texttt{Je resterende teamwoorden zijn: beha, das, narcis.} \\
\texttt{Resterende woorden voor je tegenstander zijn: mosterd, dak.} \\
\texttt{Resterende afleidingswoorden zijn: plafond, tuberculose, grot, kamer, pest.} \\
\texttt{Resterende killwoorden zijn: ontvoering.} \\
            }
        }
    }
    & \\ \\

    \theutterance \stepcounter{utterance}  
    & \multicolumn{4}{p{0.6\linewidth}}{
        \cellcolor[rgb]{0.9,0.9,0.9}{
            \makecell[{{p{\linewidth}}}]{
                \texttt{\tiny{[P1$\rangle$GM]}}
                \texttt{HINT: lingerie} \\
\texttt{DOELWITTEN: beha} \\
            }
        }
    }
    & & \\ \\

    \theutterance \stepcounter{utterance}  
    & & & \multicolumn{2}{p{0.3\linewidth}}{
        \cellcolor[rgb]{0.9,0.9,0.9}{
            \makecell[{{p{\linewidth}}}]{
                \texttt{\tiny{[GM$|$GM]}}
                \texttt{lingerie} \\
            }
        }
    }
    & & \\ \\

    \theutterance \stepcounter{utterance}  
    & & & \multicolumn{2}{p{0.3\linewidth}}{
        \cellcolor[rgb]{0.9,0.9,0.9}{
            \makecell[{{p{\linewidth}}}]{
                \texttt{\tiny{[GM$|$GM]}}
                ['beha']
            }
        }
    }
    & & \\ \\

    \theutterance \stepcounter{utterance}  
    & & & \multicolumn{2}{p{0.3\linewidth}}{
        \cellcolor[rgb]{0.9,0.9,0.9}{
            \makecell[{{p{\linewidth}}}]{
                \texttt{\tiny{[GM$|$GM]}}
                {'word': 'beha', 'assignment': 'team'}
            }
        }
    }
    & & \\ \\

    \theutterance \stepcounter{utterance}  
    & & \multicolumn{4}{p{0.6\linewidth}}{
        \cellcolor[rgb]{0.9,0.9,0.9}{
            \makecell[{{p{\linewidth}}}]{
                \texttt{\tiny{[GM$\rangle$P2]}}
                \texttt{Het woord suiker is correct geraden. Je beurt is daar geëindigd.} \\
\texttt{Geef nu nog een door komma's gescheiden lijst van ten minste 1 en maximaal 1 woorden uit de volgende lijst van woorden die het beste betrekking hebben op of het meest geassocieerd zijn met het woord 'lingerie'. Vergeet niet uw antwoord te beginnen met GISSING: . Opmerking: sommige woorden zijn verwijderd uit de lijst vergeleken met eerdere verzoeken.} \\
\\ 
\texttt{grot, dak, ontvoering, kamer, plafond, tuberculose, narcis, pest, das, beha} \\
            }
        }
    }
    & \\ \\

    \theutterance \stepcounter{utterance}  
    & & & \multicolumn{4}{p{0.6\linewidth}}{
        \cellcolor[rgb]{0.9,0.9,0.9}{
            \makecell[{{p{\linewidth}}}]{
                \texttt{\tiny{[GM$\langle$P2]}}
                \texttt{GISSING: beha} \\
            }
        }
    }
     \\ \\

    \theutterance \stepcounter{utterance}  
    & & & \multicolumn{2}{p{0.3\linewidth}}{
        \cellcolor[rgb]{0.9,0.9,0.9}{
            \makecell[{{p{\linewidth}}}]{
                \texttt{\tiny{[GM$|$GM]}}
                ['beha']
            }
        }
    }
    & & \\ \\

    \theutterance \stepcounter{utterance}  
    & & & \multicolumn{2}{p{0.3\linewidth}}{
        \cellcolor[rgb]{0.9,0.9,0.9}{
            \makecell[{{p{\linewidth}}}]{
                \texttt{\tiny{[GM$|$GM]}}
                {'word': 'beha', 'assignment': 'team'}
            }
        }
    }
    & & \\ \\

    \theutterance \stepcounter{utterance}  
    & & & \multicolumn{2}{p{0.3\linewidth}}{
        \cellcolor[rgb]{0.9,0.9,0.9}{
            \makecell[{{p{\linewidth}}}]{
                \texttt{\tiny{[GM$|$GM]}}
                {'word': 'beha', 'assignment': 'team'}
            }
        }
    }
    & & \\ \\

    \theutterance \stepcounter{utterance}  
    & & & \multicolumn{2}{p{0.3\linewidth}}{
        \cellcolor[rgb]{0.9,0.9,0.9}{
            \makecell[{{p{\linewidth}}}]{
                \texttt{\tiny{[GM$|$GM]}}
                {'word': 'dak', 'assignment': 'opponent'}
            }
        }
    }
    & & \\ \\

    \theutterance \stepcounter{utterance}  
    & & & \multicolumn{2}{p{0.3\linewidth}}{
        \cellcolor[rgb]{0.9,0.9,0.9}{
            \makecell[{{p{\linewidth}}}]{
                \texttt{\tiny{[GM$|$GM]}}
                {'hidden': {'team': ['das', 'narcis'], 'innocent': ['plafond', 'tuberculose', 'grot', 'kamer', 'pest'], 'opponent': [], 'assassin': ['ontvoering']}, 'revealed': {'team': {'team': ['sering', 'verduistering', 'zilver', 'duplex', 'tipi', 'suiker', 'beha'], 'innocent': ['pioenroos', 't-shirt'], 'opponent': ['hut'], 'assassin': []}, 'opponent': {'team': [], 'innocent': [], 'opponent': ['snelheidsovertreding', 'riem', 'wit', 'muur', 'petunia', 'mosterd', 'dak'], 'assassin': []}}}
            }
        }
    }
    & & \\ \\

    \theutterance \stepcounter{utterance}  
    & & \multicolumn{4}{p{0.6\linewidth}}{
        \cellcolor[rgb]{0.9,0.9,0.9}{
            \makecell[{{p{\linewidth}}}]{
                \texttt{\tiny{[P1$\langle$GM]}}
                \texttt{Het woord beha is correct geraden. De beurt van je teamgenoot is daar geëindigd.} \\
\texttt{Geef nu nog een aanwijzing die betrekking heeft op enkele van je overgebleven teamwoorden en een lijst met de bijbehorende doelwoorden. Vergeet niet je aanwijzing te beginnen met 'HINT: ', zet een nieuwe regel en begin je door komma's gescheiden lijst van doelwoorden met 'DOELWITTEN: '. Let op: sommige woorden zijn verwijderd uit de lijsten vergeleken met eerdere verzoeken.} \\
\\ 
\texttt{Je resterende teamwoorden zijn: das, narcis.} \\
\texttt{Resterende woorden voor je tegenstander zijn: dak.} \\
\texttt{Resterende afleidingswoorden zijn: plafond, tuberculose, grot, kamer, pest.} \\
\texttt{Resterende killwoorden zijn: ontvoering.} \\
            }
        }
    }
    & \\ \\

    \theutterance \stepcounter{utterance}  
    & \multicolumn{4}{p{0.6\linewidth}}{
        \cellcolor[rgb]{0.9,0.9,0.9}{
            \makecell[{{p{\linewidth}}}]{
                \texttt{\tiny{[P1$\rangle$GM]}}
                \texttt{HINT: natuur} \\
\texttt{DOELWITTEN: das, narcis} \\
            }
        }
    }
    & & \\ \\

    \theutterance \stepcounter{utterance}  
    & & & \multicolumn{2}{p{0.3\linewidth}}{
        \cellcolor[rgb]{0.9,0.9,0.9}{
            \makecell[{{p{\linewidth}}}]{
                \texttt{\tiny{[GM$|$GM]}}
                \texttt{natuur} \\
            }
        }
    }
    & & \\ \\

    \theutterance \stepcounter{utterance}  
    & & & \multicolumn{2}{p{0.3\linewidth}}{
        \cellcolor[rgb]{0.9,0.9,0.9}{
            \makecell[{{p{\linewidth}}}]{
                \texttt{\tiny{[GM$|$GM]}}
                ['das', 'narcis']
            }
        }
    }
    & & \\ \\

    \theutterance \stepcounter{utterance}  
    & & & \multicolumn{2}{p{0.3\linewidth}}{
        \cellcolor[rgb]{0.9,0.9,0.9}{
            \makecell[{{p{\linewidth}}}]{
                \texttt{\tiny{[GM$|$GM]}}
                {'word': 'das', 'assignment': 'team'}
            }
        }
    }
    & & \\ \\

    \theutterance \stepcounter{utterance}  
    & & & \multicolumn{2}{p{0.3\linewidth}}{
        \cellcolor[rgb]{0.9,0.9,0.9}{
            \makecell[{{p{\linewidth}}}]{
                \texttt{\tiny{[GM$|$GM]}}
                {'word': 'narcis', 'assignment': 'team'}
            }
        }
    }
    & & \\ \\

    \theutterance \stepcounter{utterance}  
    & & & \multicolumn{2}{p{0.3\linewidth}}{
        \cellcolor[rgb]{0.9,0.9,0.9}{
            \makecell[{{p{\linewidth}}}]{
                \texttt{\tiny{[GM$|$GM]}}
                \texttt{opponent has won} \\
            }
        }
    }
    & & \\ \\

\end{supertabular}
}

\end{document}
